% !TEX root = ../notes.tex

\documentclass[../notes.tex]{subfiles}

\begin{document}

\section{May 4}
Here we go.

\subsection{The Banach--Colmez Space}
Recall that
\[L_{K(h-1)}L_{K(h)}\mathbb S=E(L,\Gamma_{h-1})^{hG_{h,h-1}}.\]
This allows us to realize $L_{K(h-1)}L_{K(h)}\mathbb S$ as the totalization of the cosimplicial object
\[E(L)\to E(\mathrm{Mor}_{\mathrm{cts}}(G_{h,h-1},L))\to E(\mathrm{Mor}_{\mathrm{cts}}(G_{h,h-1}^2,L))\to\cdots.\]
One can compute this in certain cases.
\begin{remark}
	Fix $h=2$. Then Shimamora and Yabe computed the spectral sequence for $p\ge4$. Maholwad, Rezk, and Goerss computed the spectral sequence for $p=3$, and Beaudry computed it for $p=2$.
\end{remark}
One difficulty in this calculation is that Dieudonn\'e theory does not apply because $L$ is not perfect. However, we may to its perfect completion $\widehat L$. Thus, if we replace each $L$ in each object in our cosimplicial object with an $\widehat L$, then we get some object $E(\widehat L)^{hG^1_h}$, which we could hope to compute.

Anyway, let's return to our discussion of diamonds. We will be interested in the punctured Banach--Colmez space, which we now set out to define.
\begin{theorem}
	Fix an affinoid perfectoid $(R,R^+)$, and set $S\coloneqq\op{Spd}(R,R^+)$. Then $\op{Spd}\widehat L(S)$ classifies surjective maps $\OO(1/h)\to\OO(1/(h-1))$ of two vector bundles on $X_S$.
\end{theorem}
\begin{remark}
	Here, $\OO(1/h)$ is induced from a Dieudonn\'e module, which then is induced from a $p$-divisible group; namely, it comes from $\Gamma_h$. It has rank $h-1$.
\end{remark}
\begin{remark}
	By taking the kernel and twisting, it is equivalent to classify exact sequences
	\[0\to\OO\to\OO\left(\frac1h\right)\to\OO\left(\frac1{h-1}\right)\to0\]
	which preserve a choice of identification $\det\OO(1/h)\cong\det\OO(1/(h-1))$.
\end{remark}
We can now take our quotient. Viewing $G_h$ as the automorphisms of the Honda formal group, we receive a map $\det\colon G_h\to G_1$ given by taking the determinant of the corresponding Dieudonn\'e modules.
\begin{corollary}
	The quotient $\op{Spd}(\widehat L)/G_h^1$ classifies nontrivial exact sequences
	\[0\to\OO\to\mc E\to\OO\left(\frac1{h-1}\right)\to0,\]
	where $\mc E$ is isomorphic to $\OO(1/h)$ (but without a chosen isomorphism).
\end{corollary}
Now, there is a diamond which classifies extensions on $X_S$ of $\OO$ by $\OO(1/(h-1))$, denoted $\op{BC}(-1/(h-1))$, where $\op{BC}$ stands for Banach--Colmez. In other words, $\op{BC}(-1/(h-1))$ sends $(R,R^+)$ to the cohomology group $\mathrm H^1(X_{(R,R^+)};\OO(-1/(h-1)))$.

In the end, we will shortly be interested in the following variant.
\begin{notation}
	We define the diamond $\op{BC}(-1/(h-1))^*$ to parameterize nontrivial extensions
	\[0\to\OO\to\mc E\to\OO\left(\frac{-1}{h-1}\right)\to0.\]
	This is called the punctured Banach--Colmez space.
\end{notation}

\subsection{The Cohomology of Banach--Colmez Spaces}
In our language of diamonds, our present goal in homotopy theory corresponds to trying to compute
\[\mathrm R\Gamma\left(\op{BC}(-1/(h-1))^*;\OO\right).\]
It is useful to start by passing to the algebraic closure.
\begin{proposition}[Drinfeld's lemma] \label{prop:drinfeld}
	Fix an algebraically closed perfectoid field $C$ containing $k\coloneqq\FF_q$, and define $\mc C_q$ to be the $q$-Frobenius. Then
	\[\mathrm R\Gamma\left(\op{BC}\left(\frac{-1}{h-1}\right)^*;\OO\right)=\mathrm R\Gamma\left(\op{BC}\left(\frac{-1}{h-1}\right)^*\times_{\op{Spd}k}\op{Spd}C\right)^{\mc C_r}.\]
\end{proposition}
\begin{proof}[Sketch]
	Let's prove this where $\op{BC}(-1/(h-1))^*$ is replaced with $X=\op{Spd}\widehat L$, and we'll only check that
	\[\mathrm H^0(X\times_{\op{Spd}k}\op{Spd}C;\OO)^{h\mc C_q}=k((t^{1/p^\infty})).\]
	Here, once we've identified $\widehat L$ with $\op{Spd}k((t^{1/p^\infty}))$, the left-hand side consists of elements of the form
	\[\sum_{\alpha\in\ZZ[1/p]}\gamma_\alpha t^\alpha,\]
	where $\gamma_\alpha\in C$, and the sum is required to converge. Here, convergence in an interval $[r,R]$, where $r\le R$ are in $(0,1)$, means that $\left|\gamma_\alpha\right|\max\{r^\alpha,R^\alpha\}\to0$ as $\alpha\to\infty$.

	Now, to compute $\mathrm H^0$, we have to compute the kernel and cokernel of $\mc C_q-\id$, where $\mc C_r$ acts by $\sum_\alpha\gamma_\alpha t^\alpha\mapsto\sum_\alpha\gamma_\alpha^qt^\alpha$. Well, the kernel is simply those power series for which $\gamma_\alpha^q=\gamma_\alpha$, which means $\gamma_\alpha\in k$ always, and adding in convergence yields exactly power series in $k((t^{1/p^\infty}))$.

	It now remains to show that the cokernel vanishes. Well, choose some convergent $f=\sum_\alpha\gamma_\alpha t^\alpha$, and we need to find convergent $g=\sum_\alpha\delta_\alpha$ for which $\gamma_\alpha=\delta_\alpha^q-\delta_\alpha$. Certainly such a thing exists because $C$ is algebraically closed, and in fact, one can find $\delta_\alpha$ with $\left|\delta_\alpha\right|\le\left|\gamma_\alpha\right|^{1/q}$; for example, given any solution $\delta_\alpha$, the other solutions look like $\delta_\alpha+c$ for $c\in\FF_q$, and we just need to choose $c$ correctly. Given this bound, one finds that $f_\alpha$ converging on an interval $[r,R]$ implies that $g$ converges on $\left[r^{1/q},R^{1/q}\right]$.
\end{proof}
For convenience, we fix the following choice of algebraic closure.
\begin{notation}
	Let $C$ be the completion of the algebraic closure of $W(k)[1/p]$. Note that $C$ has characteristic zero. Then for any $t$, we define
	\[\mathrm R\Gamma(\op{BC}(-1/t)\times_{\op{Spd}k}\op{Spd}C;\OO^\flat)=\mathrm R\Gamma\left(\op{BC}(-1/t)\times_{\op{Spd}k}\op{Spd}C^\flat;\OO\right)\]
\end{notation}
Let's now compute with $\op{BC}(-1/t)$.
\begin{proposition} \label{prob:bc-as-quotient}
	Fix a field extension $F/\QQ_p$ of degree $t$, which embed into $C$. Then there is an isomorphism
	\[\op{BC}(-1/t)\times_{\op{Spd}k}\op{Spd}C\cong\AA^{1,\diamond}_C/F\]
	of $\QQ_p$-linear diamonds.
\end{proposition}
\begin{proof}[Sketch]
	A test object $S^\flat\op{Spd}D$ provides the the data of an untilt $(R,R^+)$, meaning that $(R^\flat,R^{\flat+})=S^\flat$, along with a point $i\colon S\to X_{S^\flat}$. Now, by passing to characteristic zero, there is an exact sequence
	\[0\to\OO_{X_{S^\flat,F}}\to\OO_{X_{S^\flat,F}}\to i_*\OO_S\to0\]
	of quasicoherent sheaves. Now, let $\pi\colon X_{S^\flat,F}\to X_{S^\flat}$ be the \'etale covering from the base-change $\QQ_p\to F$. Pushing forward this exact sequence along $\pi$, one receives a short exact sequence
	\[0\to\OO_{X_{S^\flat}}(-1/t)\to\OO_{X_{S^\flat}}\otimes_{\QQ_p}F\to(\pi i)_*\OO_S\to0.\]
	We now take the long exact sequence in cohomology. One has $\mathrm H^0(X_{S^\flat};\OO(-1))=0$ and $\mathrm H^0(X_{S^\flat};\OO)=\QQ_p$ and $\mathrm H^0(X_{S^\flat};i_*\OO_S)=\OO$ (where $\OO$ is the functor sending $S\mapsto R$), so we have an exact sequence
	\[0\to \QQ_p\otimes_{\QQ_p}F\to\AA^1\to\mathrm H^1(\OO_{X_{S^\flat}}(-1/t))\to0.\]
	Thus, we see that $\op{BC}(-1/t)$ is the same as the quotient of $\AA^1$ by $F$ on test objects, so we are done.
\end{proof}

\end{document}